\section{Modern webalkalmazások fejlesztési folyamata}

A modern webalkalmazások fejlesztése általában a szoftverfejlesztési életciklus, vagyis az SDLC (Software Development Life Cycle) lépéseire épül. Az SDLC célja, hogy a fejlesztés előre meghatározott, átlátható lépések mentén haladjon: a csapat először meghatározza a követelményeket, majd megtervezi, megvalósítja, teszteli, telepíti és hosszú távon karbantartja az alkalmazást \cite{ibmSdlc}. Webalkalmazások esetében ez a folyamat különösen fontos, mivel a rendszer több rétegből áll: frontendből, backendből, adatbázisból, API-kból, infrastruktúrából és sok esetben külső szolgáltatásokból is \cite{mdnClientServerOverview,mdnServerSideIntro}.

A fejlesztési folyamat első szakasza a követelmények összegyűjtése és pontosítása \cite{isoIecIeee29148Requirements}. Ebben a fázisban a csapat meghatározza, hogy az alkalmazás milyen problémát old meg, kik lesznek a felhasználói, milyen funkciókra van szükség, és milyen üzleti vagy technikai korlátokat kell figyelembe venni. Ide tartozhatnak a funkcionális követelmények, például a felhasználói regisztráció, keresés vagy adminisztrációs felület, valamint a nem funkcionális követelmények is, például a teljesítmény, biztonság, skálázhatóság, hozzáférhetőség és reszponzív működés \cite{isoIec25010ProductQuality,w3cWaiAccessibilityIntro}.

A következő lépés a tervezés, amely során kialakul az alkalmazás architektúrája és felhasználói felülete \cite{isoIecIeee42010Architecture,nielsenNorman1998ux}. A UX/UI tervezés meghatározza a felhasználói útvonalakat, a képernyők felépítését és az interakciók logikáját, míg a technikai tervezés a frontend és backend feladatok szétválasztására, az API-k kialakítására, az adatmodellre és az infrastruktúrára koncentrál \cite{mdnClientServerOverview}. Egy jól átgondolt tervezési szakasz csökkenti a későbbi félreértéseket, és segít abban, hogy a fejlesztők egységes képet kapjanak a megvalósítandó rendszerről.

Az implementáció során a fejlesztők elkészítik az alkalmazás működő részeit. A frontend fejlesztés a böngészőben futó felületre, a backend fejlesztés pedig a szerveroldali logikára, adatkezelésre, autentikációra, autorizációra és API-kra koncentrál \cite{mdnWebStandardsModel,mdnServerSideIntro}. A modern webfejlesztésben gyakori, hogy a fejlesztés iteratív módon történik: a csapat kisebb funkciókat valósít meg, ezeket folyamatosan teszteli, majd visszajelzések alapján továbbfejleszti \cite{agileManifesto2001}.

A tesztelés a fejlesztési folyamat egyik kulcsfontosságú része. Nemcsak a kész alkalmazás ellenőrzését jelenti, hanem a teljes fejlesztés során jelen van \cite{ibmShiftLeftTesting}. A manuális tesztek mellett egyre nagyobb szerepet kapnak az automatizált tesztek, például az egységtesztek, integrációs tesztek és end-to-end tesztek \cite{mdnAutomatedTesting}. Ezek segítenek kiszűrni a regressziós hibákat is, amikor egy új fejlesztés vagy hibajavítás elront egy korábban már működő funkciót \cite{istqbGlossaryRegressionTesting}.

A telepítés és üzemeltetés szakaszában az alkalmazás kikerül a felhasználók számára elérhető környezetbe. Modern fejlesztési környezetekben ezt gyakran CI/CD folyamatok támogatják, amelyek automatizálják a buildelést, tesztelést és telepítést \cite{awsCicd}. A DevOps szemlélet célja, hogy a fejlesztés és az üzemeltetés ne különálló, egymástól elszigetelt tevékenység legyen, hanem egységes, együttműködő folyamat \cite{atlassianDevOps}.

Az SDLC utolsó, de folyamatosan visszatérő része a karbantartás és továbbfejlesztés \cite{isoIecIeee14764Maintenance}. Egy webalkalmazás a kiadás után sem tekinthető befejezettnek: hibajavításokra, biztonsági frissítésekre, teljesítményoptimalizálásra és új funkciók bevezetésére is szükség lehet \cite{nistSp800218Ssdf}. Emiatt a modern webalkalmazás-fejlesztés inkább ciklikus folyamatként értelmezhető, ahol a tervezés, fejlesztés, tesztelés és üzemeltetés folyamatosan ismétlődik.
